%%%%%%%%%%%%%%%%%%%%%%%%%%%%%%%%
% Chapter 5. Conclusion
%%%%%%%%%%%%%%%%%%%%%%%%%%%%%%%%

\section{결론}\label{sec:conclusion}

본 논문은 사전학습 임베딩 행렬의 항별 통계에 맞춘 등방 가우시안에서 추첨한 RSP가 학습 없이도 영향을 줄 수 있는 이유를 transformer 수준에서 분석했다. 분석의 핵심은 네 가지다. exact attention decomposition, 초반 exploration의 lower envelope(§\ref{sec:improve}), 후반 감쇠의 upper envelope(§\ref{sec:converge}), 그리고 독립 재추첨과 별도의 과제 측 성공 가정이 더해질 때에만 Pass@$N$ 증가로 이어지는 branch-opening 논리(§\ref{sec:multipath})다. §\ref{sec:experimental_analysis}의 실험은 이 구조와 맞물린다. 학습 없는 RSP는 경쟁적인 성능을 유지하고, 초반 토큰 분포는 더 다양해진 뒤 후반에 안정화되며, 같은 RSP를 모든 sample에 공유하면 누적 이득이 사라진다. RSP는 순위를 \emph{바꾸는} 섭동이고 temperature는 순위를 \emph{보존}하는 섭동이므로 둘은 직교적이며, 이는 GRPO \citep{shao2024deepseekmath}의 entropy collapse를 DAPO \citep{yu2025dapo}, CE-GPPO \citep{su2025cegppo} 너머에서 샘플링 단계로 보완할 가능성을 시사한다. RSP를 활용한 학습 시점 비교, 수학 추론을 넘어선 도메인 확장, RSP 길이·주입 위치의 원리적 결정은 향후 과제다.
